
\documentclass{article}
\usepackage{colortbl}
\usepackage{makecell}
\usepackage{multirow}
\usepackage{supertabular}

\begin{document}

\newcounter{utterance}

\centering \large Interaction Transcript for game `matchit\_ascii', experiment `similar\_grid\_2', episode 0 with nemotron{-}3{-}ultra.
\vspace{24pt}

{ \footnotesize  \setcounter{utterance}{1}
\setlength{\tabcolsep}{0pt}
\begin{supertabular}{c@{$\;$}|p{.15\linewidth}@{}p{.15\linewidth}p{.15\linewidth}p{.15\linewidth}p{.15\linewidth}p{.15\linewidth}}
    \# & $\;$A & \multicolumn{4}{c}{Game Master} & $\;\:$B\\
    \hline

    \theutterance \stepcounter{utterance}  
    & & \multicolumn{4}{p{0.6\linewidth}}{
        \cellcolor[rgb]{0.9,0.9,0.9}{
            \makecell[{{p{\linewidth}}}]{
                \texttt{\tiny{[P1$\langle$GM]}}
                \texttt{Osallistut yhteistyöhön perustuvaan arvauspeliin.} \\
\texttt{Tavoitteena on selvittää, ovatko tämä ruudukko ja toinen kuva, jonka vain minä näen, samat.} \\
\texttt{Kuvaile ensin ruudukkosi. Sitten annan oman kuvaukseni ja voimme esittää toisillemme kysymyksiä ruudukoista selvittääksemme, ovatko ne samat. Aloita nyt lyhyt ruudukko{-}kuvauksesi merkillä "KUVAUS:" ja sen jälkeen kuvaus. Älä lisää mitään muuta. Tämä on sinun ruudukkosi:} \\
\texttt{▢ X X X X} \\
\texttt{▢ X X X X} \\
\texttt{▢ X X X X} \\
\texttt{▢ X X X X} \\
\texttt{▢ ▢ ▢ ▢ ▢} \\
            }
        }
    }
    & \\ \\

    \theutterance \stepcounter{utterance}  
    & \multicolumn{4}{p{0.6\linewidth}}{
        \cellcolor[rgb]{0.9,0.9,0.9}{
            \makecell[{{p{\linewidth}}}]{
                \texttt{\tiny{[P1$\rangle$GM]}}
                \texttt{KUVAUS: Ruudukko on 5x4{-}kokoinen, mutta viidennellä rivillä on 5 ruutua. Ensimmäiset neljä riviä: ensimmäinen ruutu tyhjä, seuraavat neljä on X{-}merkatut. Viides rivi: kaikki viisi ruutua tyhjiä.} \\
            }
        }
    }
    & & \\ \\

    \theutterance \stepcounter{utterance}  
    & & & \multicolumn{2}{p{0.3\linewidth}}{
        \cellcolor[rgb]{0.9,0.9,0.9}{
            \makecell[{{p{\linewidth}}}]{
                \texttt{\tiny{[GM$|$GM]}}
                \texttt{continue} \\
            }
        }
    }
    & & \\ \\

    \theutterance \stepcounter{utterance}  
    & & \multicolumn{4}{p{0.6\linewidth}}{
        \cellcolor[rgb]{0.9,0.9,0.9}{
            \makecell[{{p{\linewidth}}}]{
                \texttt{\tiny{[GM$\rangle$P2]}}
                \texttt{Osallistut yhteistyöhön perustuvaan arvauspeliin.} \\
\texttt{Tavoitteena on selvittää, ovatko tämä ruudukko ja toinen kuva, jonka vain minä näen, samat.} \\
\texttt{Kuvaile ensin ruudukkosi. Sitten annan oman kuvaukseni ja voimme esittää toisillemme kysymyksiä ruudukoista selvittääksemme, ovatko ne samat. Aloita nyt lyhyt ruudukko{-}kuvauksesi merkillä "KUVAUS:" ja sen jälkeen kuvaus. Älä lisää mitään muuta. Tämä on sinun ruudukkosi:} \\
\texttt{X X X X X} \\
\texttt{▢ X X X ▢} \\
\texttt{▢ X X X X} \\
\texttt{▢ X X X X} \\
\texttt{▢ ▢ ▢ ▢ ▢} \\
            }
        }
    }
    & \\ \\

    \theutterance \stepcounter{utterance}  
    & & & \multicolumn{4}{p{0.6\linewidth}}{
        \cellcolor[rgb]{0.9,0.9,0.9}{
            \makecell[{{p{\linewidth}}}]{
                \texttt{\tiny{[GM$\langle$P2]}}
                \texttt{KUVAUS: 5x5{-}ruudukko. Rivit ylä alaspäin: 1. täysin täytetty (5 X:ää). 2. tyhjä, 3 X:ää, tyhjä. 3. tyhjä, 4 X:ää. 4. tyhjä, 4 X:ää. 5. täysin tyhjä (5 tyhjää).} \\
            }
        }
    }
     \\ \\

    \theutterance \stepcounter{utterance}  
    & & & \multicolumn{2}{p{0.3\linewidth}}{
        \cellcolor[rgb]{0.9,0.9,0.9}{
            \makecell[{{p{\linewidth}}}]{
                \texttt{\tiny{[GM$|$GM]}}
                \texttt{continue} \\
            }
        }
    }
    & & \\ \\

    \theutterance \stepcounter{utterance}  
    & & \multicolumn{4}{p{0.6\linewidth}}{
        \cellcolor[rgb]{0.9,0.9,0.9}{
            \makecell[{{p{\linewidth}}}]{
                \texttt{\tiny{[GM$\rangle$P2]}}
                \texttt{Tämä on KUVAUKSENIKUVAUS: Ruudukko on 5x4{-}kokoinen, mutta viidennellä rivillä on 5 ruutua. Ensimmäiset neljä riviä: ensimmäinen ruutu tyhjä, seuraavat neljä on X{-}merkatut. Viides rivi: kaikki viisi ruutua tyhjiä.} \\
\texttt{Esitä nyt kysymys selvittääksesi uusia näkökohtia ruudukostani, jotka saattavat poiketa sinun ruudukostasi. Aloita "KYSYMYS:" äläkä lisää mitään muuta.} \\
            }
        }
    }
    & \\ \\

    \theutterance \stepcounter{utterance}  
    & & & \multicolumn{4}{p{0.6\linewidth}}{
        \cellcolor[rgb]{0.9,0.9,0.9}{
            \makecell[{{p{\linewidth}}}]{
                \texttt{\tiny{[GM$\langle$P2]}}
                \texttt{KYSYMYS: Onko sinun ruudukossasi tarkalleen viisi riviä (eli onko se 5x5{-}ruudukko)?} \\
            }
        }
    }
     \\ \\

    \theutterance \stepcounter{utterance}  
    & & & \multicolumn{2}{p{0.3\linewidth}}{
        \cellcolor[rgb]{0.9,0.9,0.9}{
            \makecell[{{p{\linewidth}}}]{
                \texttt{\tiny{[GM$|$GM]}}
                \texttt{continue} \\
            }
        }
    }
    & & \\ \\

    \theutterance \stepcounter{utterance}  
    & & \multicolumn{4}{p{0.6\linewidth}}{
        \cellcolor[rgb]{0.9,0.9,0.9}{
            \makecell[{{p{\linewidth}}}]{
                \texttt{\tiny{[P1$\langle$GM]}}
                \texttt{Tämä on KUVAUKSENIKUVAUS: 5x5{-}ruudukko. Rivit ylä alaspäin: 1. täysin täytetty (5 X:ää). 2. tyhjä, 3 X:ää, tyhjä. 3. tyhjä, 4 X:ää. 4. tyhjä, 4 X:ää. 5. täysin tyhjä (5 tyhjää).} \\
\texttt{KYSYMYS: Onko sinun ruudukossasi tarkalleen viisi riviä (eli onko se 5x5{-}ruudukko)?Aloita vastauksesi merkillä "VASTAUS:" äläkä lisää mitään muuta.} \\
            }
        }
    }
    & \\ \\

    \theutterance \stepcounter{utterance}  
    & \multicolumn{4}{p{0.6\linewidth}}{
        \cellcolor[rgb]{0.9,0.9,0.9}{
            \makecell[{{p{\linewidth}}}]{
                \texttt{\tiny{[P1$\rangle$GM]}}
                \texttt{VASTAUS: Kyllä, ruudukossani on viisi riviä, mutta se ei ole 5x5{-}ruudukko, sillä viidellä rivillä on viisi ruutua, kun ensimmäisillä neljällä on neljä.} \\
            }
        }
    }
    & & \\ \\

    \theutterance \stepcounter{utterance}  
    & & & \multicolumn{2}{p{0.3\linewidth}}{
        \cellcolor[rgb]{0.9,0.9,0.9}{
            \makecell[{{p{\linewidth}}}]{
                \texttt{\tiny{[GM$|$GM]}}
                \texttt{continue} \\
            }
        }
    }
    & & \\ \\

    \theutterance \stepcounter{utterance}  
    & & \multicolumn{4}{p{0.6\linewidth}}{
        \cellcolor[rgb]{0.9,0.9,0.9}{
            \makecell[{{p{\linewidth}}}]{
                \texttt{\tiny{[P1$\langle$GM]}}
                \texttt{Esitä nyt kysymys selvittääksesi uusia näkökohtia ruudukostani, jotka saattavat poiketa sinun ruudukostasi. Aloita "KYSYMYS:" äläkä lisää mitään muuta.} \\
            }
        }
    }
    & \\ \\

    \theutterance \stepcounter{utterance}  
    & \multicolumn{4}{p{0.6\linewidth}}{
        \cellcolor[rgb]{0.9,0.9,0.9}{
            \makecell[{{p{\linewidth}}}]{
                \texttt{\tiny{[P1$\rangle$GM]}}
                \texttt{KYSYMYS: Ensimmäisellä rivillä on viisi X:ää peräkkäin. Onko sinun ruudukossasi minkään rivin alkuruutu X?} \\
            }
        }
    }
    & & \\ \\

    \theutterance \stepcounter{utterance}  
    & & & \multicolumn{2}{p{0.3\linewidth}}{
        \cellcolor[rgb]{0.9,0.9,0.9}{
            \makecell[{{p{\linewidth}}}]{
                \texttt{\tiny{[GM$|$GM]}}
                \texttt{continue} \\
            }
        }
    }
    & & \\ \\

    \theutterance \stepcounter{utterance}  
    & & \multicolumn{4}{p{0.6\linewidth}}{
        \cellcolor[rgb]{0.9,0.9,0.9}{
            \makecell[{{p{\linewidth}}}]{
                \texttt{\tiny{[GM$\rangle$P2]}}
                \texttt{VASTAUS: Kyllä, ruudukossani on viisi riviä, mutta se ei ole 5x5{-}ruudukko, sillä viidellä rivillä on viisi ruutua, kun ensimmäisillä neljällä on neljä.} \\
\texttt{KYSYMYS: Ensimmäisellä rivillä on viisi X:ää peräkkäin. Onko sinun ruudukossasi minkään rivin alkuruutu X?Aloita vastauksesi merkillä "VASTAUS:" äläkä lisää mitään muuta.} \\
            }
        }
    }
    & \\ \\

    \theutterance \stepcounter{utterance}  
    & & & \multicolumn{4}{p{0.6\linewidth}}{
        \cellcolor[rgb]{0.9,0.9,0.9}{
            \makecell[{{p{\linewidth}}}]{
                \texttt{\tiny{[GM$\langle$P2]}}
                \texttt{VASTAUS: Kyllä, ensimmäinen rivi alkaa X:llä (ja koko rivi on X:itä).} \\
            }
        }
    }
     \\ \\

    \theutterance \stepcounter{utterance}  
    & & & \multicolumn{2}{p{0.3\linewidth}}{
        \cellcolor[rgb]{0.9,0.9,0.9}{
            \makecell[{{p{\linewidth}}}]{
                \texttt{\tiny{[GM$|$GM]}}
                \texttt{continue} \\
            }
        }
    }
    & & \\ \\

    \theutterance \stepcounter{utterance}  
    & & \multicolumn{4}{p{0.6\linewidth}}{
        \cellcolor[rgb]{0.9,0.9,0.9}{
            \makecell[{{p{\linewidth}}}]{
                \texttt{\tiny{[GM$\rangle$P2]}}
                \texttt{Esitä nyt kysymys selvittääksesi uusia näkökohtia ruudukostani, jotka saattavat poiketa sinun ruudukostasi. Aloita "KYSYMYS:" äläkä lisää mitään muuta.} \\
            }
        }
    }
    & \\ \\

    \theutterance \stepcounter{utterance}  
    & & & \multicolumn{4}{p{0.6\linewidth}}{
        \cellcolor[rgb]{0.9,0.9,0.9}{
            \makecell[{{p{\linewidth}}}]{
                \texttt{\tiny{[GM$\langle$P2]}}
                \texttt{KYSYMYS: Onko toisen rivin ensimmäinen ruutu tyhjä?} \\
            }
        }
    }
     \\ \\

    \theutterance \stepcounter{utterance}  
    & & & \multicolumn{2}{p{0.3\linewidth}}{
        \cellcolor[rgb]{0.9,0.9,0.9}{
            \makecell[{{p{\linewidth}}}]{
                \texttt{\tiny{[GM$|$GM]}}
                \texttt{continue} \\
            }
        }
    }
    & & \\ \\

    \theutterance \stepcounter{utterance}  
    & & \multicolumn{4}{p{0.6\linewidth}}{
        \cellcolor[rgb]{0.9,0.9,0.9}{
            \makecell[{{p{\linewidth}}}]{
                \texttt{\tiny{[P1$\langle$GM]}}
                \texttt{VASTAUS: Kyllä, ensimmäinen rivi alkaa X:llä (ja koko rivi on X:itä).} \\
\texttt{KYSYMYS: Onko toisen rivin ensimmäinen ruutu tyhjä?Aloita vastauksesi merkillä "VASTAUS:" äläkä lisää mitään muuta.} \\
            }
        }
    }
    & \\ \\

    \theutterance \stepcounter{utterance}  
    & \multicolumn{4}{p{0.6\linewidth}}{
        \cellcolor[rgb]{0.9,0.9,0.9}{
            \makecell[{{p{\linewidth}}}]{
                \texttt{\tiny{[P1$\rangle$GM]}}
                \texttt{VASTAUS: Kyllä, toisen rivin ensimmäinen ruutu on tyhjä.} \\
            }
        }
    }
    & & \\ \\

    \theutterance \stepcounter{utterance}  
    & & & \multicolumn{2}{p{0.3\linewidth}}{
        \cellcolor[rgb]{0.9,0.9,0.9}{
            \makecell[{{p{\linewidth}}}]{
                \texttt{\tiny{[GM$|$GM]}}
                \texttt{continue} \\
            }
        }
    }
    & & \\ \\

    \theutterance \stepcounter{utterance}  
    & & \multicolumn{4}{p{0.6\linewidth}}{
        \cellcolor[rgb]{0.9,0.9,0.9}{
            \makecell[{{p{\linewidth}}}]{
                \texttt{\tiny{[P1$\langle$GM]}}
                \texttt{Esitä nyt kysymys selvittääksesi uusia näkökohtia ruudukostani, jotka saattavat poiketa sinun ruudukostasi. Aloita "KYSYMYS:" äläkä lisää mitään muuta.} \\
            }
        }
    }
    & \\ \\

    \theutterance \stepcounter{utterance}  
    & \multicolumn{4}{p{0.6\linewidth}}{
        \cellcolor[rgb]{0.9,0.9,0.9}{
            \makecell[{{p{\linewidth}}}]{
                \texttt{\tiny{[P1$\rangle$GM]}}
                \texttt{KYSYMYS: Toisella rivillä on kuvauksesi mukaan "tyhjä, 3 X:ää, tyhjä". Ovatko kolmannen ja neljännen rivin rakenteet identtiset toisen rivin kanssa (eli tyhjä, 3 X:ää, tyhjä), vai ovatko ne "tyhjä, 4 X:ää" ilman lopussa olevaa tyhjää ruutua?} \\
            }
        }
    }
    & & \\ \\

    \theutterance \stepcounter{utterance}  
    & & & \multicolumn{2}{p{0.3\linewidth}}{
        \cellcolor[rgb]{0.9,0.9,0.9}{
            \makecell[{{p{\linewidth}}}]{
                \texttt{\tiny{[GM$|$GM]}}
                \texttt{continue} \\
            }
        }
    }
    & & \\ \\

    \theutterance \stepcounter{utterance}  
    & & \multicolumn{4}{p{0.6\linewidth}}{
        \cellcolor[rgb]{0.9,0.9,0.9}{
            \makecell[{{p{\linewidth}}}]{
                \texttt{\tiny{[GM$\rangle$P2]}}
                \texttt{VASTAUS: Kyllä, toisen rivin ensimmäinen ruutu on tyhjä.} \\
\texttt{KYSYMYS: Toisella rivillä on kuvauksesi mukaan "tyhjä, 3 X:ää, tyhjä". Ovatko kolmannen ja neljännen rivin rakenteet identtiset toisen rivin kanssa (eli tyhjä, 3 X:ää, tyhjä), vai ovatko ne "tyhjä, 4 X:ää" ilman lopussa olevaa tyhjää ruutua?Aloita vastauksesi merkillä "VASTAUS:" äläkä lisää mitään muuta.} \\
            }
        }
    }
    & \\ \\

    \theutterance \stepcounter{utterance}  
    & & & \multicolumn{4}{p{0.6\linewidth}}{
        \cellcolor[rgb]{0.9,0.9,0.9}{
            \makecell[{{p{\linewidth}}}]{
                \texttt{\tiny{[GM$\langle$P2]}}
                \texttt{VASTAUS: Kolmannen ja neljännen rivin rakenne on "tyhjä, 4 X:ää" (ilman lopussa olevaa tyhjää ruutua). Ne eroavat toisesta rivistä.} \\
            }
        }
    }
     \\ \\

    \theutterance \stepcounter{utterance}  
    & & & \multicolumn{2}{p{0.3\linewidth}}{
        \cellcolor[rgb]{0.9,0.9,0.9}{
            \makecell[{{p{\linewidth}}}]{
                \texttt{\tiny{[GM$|$GM]}}
                \texttt{continue} \\
            }
        }
    }
    & & \\ \\

    \theutterance \stepcounter{utterance}  
    & & \multicolumn{4}{p{0.6\linewidth}}{
        \cellcolor[rgb]{0.9,0.9,0.9}{
            \makecell[{{p{\linewidth}}}]{
                \texttt{\tiny{[GM$\rangle$P2]}}
                \texttt{Esitä nyt kysymys selvittääksesi uusia näkökohtia ruudukostani, jotka saattavat poiketa sinun ruudukostasi. Aloita "KYSYMYS:" äläkä lisää mitään muuta.} \\
            }
        }
    }
    & \\ \\

    \theutterance \stepcounter{utterance}  
    & & & \multicolumn{4}{p{0.6\linewidth}}{
        \cellcolor[rgb]{0.9,0.9,0.9}{
            \makecell[{{p{\linewidth}}}]{
                \texttt{\tiny{[GM$\langle$P2]}}
                \texttt{KYSYMYS: Ovatko kolmannen ja neljännen rivin ruudut viidessä sarakkeessa (eli rivin neljännen X:n jälkeen) tyhjiä?} \\
            }
        }
    }
     \\ \\

    \theutterance \stepcounter{utterance}  
    & & & \multicolumn{2}{p{0.3\linewidth}}{
        \cellcolor[rgb]{0.9,0.9,0.9}{
            \makecell[{{p{\linewidth}}}]{
                \texttt{\tiny{[GM$|$GM]}}
                \texttt{continue} \\
            }
        }
    }
    & & \\ \\

    \theutterance \stepcounter{utterance}  
    & & \multicolumn{4}{p{0.6\linewidth}}{
        \cellcolor[rgb]{0.9,0.9,0.9}{
            \makecell[{{p{\linewidth}}}]{
                \texttt{\tiny{[P1$\langle$GM]}}
                \texttt{VASTAUS: Kolmannen ja neljännen rivin rakenne on "tyhjä, 4 X:ää" (ilman lopussa olevaa tyhjää ruutua). Ne eroavat toisesta rivistä.} \\
\texttt{KYSYMYS: Ovatko kolmannen ja neljännen rivin ruudut viidessä sarakkeessa (eli rivin neljännen X:n jälkeen) tyhjiä?Aloita vastauksesi merkillä "VASTAUS:" äläkä lisää mitään muuta.} \\
            }
        }
    }
    & \\ \\

    \theutterance \stepcounter{utterance}  
    & \multicolumn{4}{p{0.6\linewidth}}{
        \cellcolor[rgb]{0.9,0.9,0.9}{
            \makecell[{{p{\linewidth}}}]{
                \texttt{\tiny{[P1$\rangle$GM]}}
                \texttt{VASTAUS: Kyllä, kolmannella ja neljällä rivillä viides ruutu (sarake 5) on tyhjä, sillä nillä on vain neljä ruutua (1 tyhjä + 4 X:ää). Viidellä rivillä viides ruutu on myös tyhjä.} \\
            }
        }
    }
    & & \\ \\

    \theutterance \stepcounter{utterance}  
    & & & \multicolumn{2}{p{0.3\linewidth}}{
        \cellcolor[rgb]{0.9,0.9,0.9}{
            \makecell[{{p{\linewidth}}}]{
                \texttt{\tiny{[GM$|$GM]}}
                \texttt{continue} \\
            }
        }
    }
    & & \\ \\

    \theutterance \stepcounter{utterance}  
    & & \multicolumn{4}{p{0.6\linewidth}}{
        \cellcolor[rgb]{0.9,0.9,0.9}{
            \makecell[{{p{\linewidth}}}]{
                \texttt{\tiny{[P1$\langle$GM]}}
                \texttt{Esitä nyt kysymys selvittääksesi uusia näkökohtia ruudukostani, jotka saattavat poiketa sinun ruudukostasi. Aloita "KYSYMYS:" äläkä lisää mitään muuta.} \\
            }
        }
    }
    & \\ \\

    \theutterance \stepcounter{utterance}  
    & \multicolumn{4}{p{0.6\linewidth}}{
        \cellcolor[rgb]{0.9,0.9,0.9}{
            \makecell[{{p{\linewidth}}}]{
                \texttt{\tiny{[P1$\rangle$GM]}}
                \texttt{KYSYMYS: Ensimmäisellä rivillä on viisi X:ää. Miten first{-}sarakkeen ruudut voivat olla samat, jos minulla ensimmäinen rivi alkaa X:llä ja sinulla tyhjällä? Oletatko varma, että ruudut ovat solumuotoisia ja rivit tasalla?} \\
            }
        }
    }
    & & \\ \\

    \theutterance \stepcounter{utterance}  
    & & & \multicolumn{2}{p{0.3\linewidth}}{
        \cellcolor[rgb]{0.9,0.9,0.9}{
            \makecell[{{p{\linewidth}}}]{
                \texttt{\tiny{[GM$|$GM]}}
                \texttt{continue} \\
            }
        }
    }
    & & \\ \\

    \theutterance \stepcounter{utterance}  
    & & \multicolumn{4}{p{0.6\linewidth}}{
        \cellcolor[rgb]{0.9,0.9,0.9}{
            \makecell[{{p{\linewidth}}}]{
                \texttt{\tiny{[GM$\rangle$P2]}}
                \texttt{VASTAUS: Kyllä, kolmannella ja neljällä rivillä viides ruutu (sarake 5) on tyhjä, sillä nillä on vain neljä ruutua (1 tyhjä + 4 X:ää). Viidellä rivillä viides ruutu on myös tyhjä.} \\
\texttt{KYSYMYS: Ensimmäisellä rivillä on viisi X:ää. Miten first{-}sarakkeen ruudut voivat olla samat, jos minulla ensimmäinen rivi alkaa X:llä ja sinulla tyhjällä? Oletatko varma, että ruudut ovat solumuotoisia ja rivit tasalla?Aloita vastauksesi merkillä "VASTAUS:" äläkä lisää mitään muuta.} \\
            }
        }
    }
    & \\ \\

    \theutterance \stepcounter{utterance}  
    & & & \multicolumn{4}{p{0.6\linewidth}}{
        \cellcolor[rgb]{0.9,0.9,0.9}{
            \makecell[{{p{\linewidth}}}]{
                \texttt{\tiny{[GM$\langle$P2]}}
                \texttt{VASTAUS: Ruudut eivät ole samat. Kysyin aiemmin, onko sinun toisen rivin ensimmäinen ruutu tyhjä, ja vastasit kyllä. Minun toisen rivin ensimmäinen ruutu on tyhjä, mutta minun *ensimmäisen* rivin ensimmäinen ruutu on X. Ensimmäinen rivi erottää ruudut toisistaan selkeästi.} \\
            }
        }
    }
     \\ \\

    \theutterance \stepcounter{utterance}  
    & & & \multicolumn{2}{p{0.3\linewidth}}{
        \cellcolor[rgb]{0.9,0.9,0.9}{
            \makecell[{{p{\linewidth}}}]{
                \texttt{\tiny{[GM$|$GM]}}
                \texttt{continue} \\
            }
        }
    }
    & & \\ \\

    \theutterance \stepcounter{utterance}  
    & & \multicolumn{4}{p{0.6\linewidth}}{
        \cellcolor[rgb]{0.9,0.9,0.9}{
            \makecell[{{p{\linewidth}}}]{
                \texttt{\tiny{[GM$\rangle$P2]}}
                \texttt{Tee nyt päätös. Mitä mieltä olet: ovatko ruudukkosi ja toinen kuvattu ruudukko samat? Kirjoita "PÄÄTÖS: sama ruudukko", jos uskot niiden olevan sama ruudukko, tai "PÄÄTÖS: eri ruudukot", jos uskot niiden olevan eri ruudukot. Älä lisää mitään muuta.} \\
            }
        }
    }
    & \\ \\

    \theutterance \stepcounter{utterance}  
    & & & \multicolumn{4}{p{0.6\linewidth}}{
        \cellcolor[rgb]{0.9,0.9,0.9}{
            \makecell[{{p{\linewidth}}}]{
                \texttt{\tiny{[GM$\langle$P2]}}
                \texttt{PÄÄTÖS: eri ruudukot} \\
            }
        }
    }
     \\ \\

    \theutterance \stepcounter{utterance}  
    & & & \multicolumn{2}{p{0.3\linewidth}}{
        \cellcolor[rgb]{0.9,0.9,0.9}{
            \makecell[{{p{\linewidth}}}]{
                \texttt{\tiny{[GM$|$GM]}}
                \texttt{continue} \\
            }
        }
    }
    & & \\ \\

    \theutterance \stepcounter{utterance}  
    & & & \multicolumn{2}{p{0.3\linewidth}}{
        \cellcolor[rgb]{0.9,0.9,0.9}{
            \makecell[{{p{\linewidth}}}]{
                \texttt{\tiny{[GM$|$GM]}}
                \texttt{success} \\
            }
        }
    }
    & & \\ \\

    \theutterance \stepcounter{utterance}  
    & & \multicolumn{4}{p{0.6\linewidth}}{
        \cellcolor[rgb]{0.9,0.9,0.9}{
            \makecell[{{p{\linewidth}}}]{
                \texttt{\tiny{[P1$\langle$GM]}}
                \texttt{VASTAUS: Ruudut eivät ole samat. Kysyin aiemmin, onko sinun toisen rivin ensimmäinen ruutu tyhjä, ja vastasit kyllä. Minun toisen rivin ensimmäinen ruutu on tyhjä, mutta minun *ensimmäisen* rivin ensimmäinen ruutu on X. Ensimmäinen rivi erottää ruudut toisistaan selkeästi.} \\
\texttt{Tee nyt päätös. Mitä mieltä olet: ovatko ruudukkosi ja toinen kuvattu ruudukko samat? Kirjoita "PÄÄTÖS: sama ruudukko", jos uskot niiden olevan sama ruudukko, tai "PÄÄTÖS: eri ruudukot", jos uskot niiden olevan eri ruudukot. Älä lisää mitään muuta.} \\
            }
        }
    }
    & \\ \\

    \theutterance \stepcounter{utterance}  
    & \multicolumn{4}{p{0.6\linewidth}}{
        \cellcolor[rgb]{0.9,0.9,0.9}{
            \makecell[{{p{\linewidth}}}]{
                \texttt{\tiny{[P1$\rangle$GM]}}
                \texttt{PÄÄTÖS: eri ruudukot} \\
            }
        }
    }
    & & \\ \\

    \theutterance \stepcounter{utterance}  
    & & & \multicolumn{2}{p{0.3\linewidth}}{
        \cellcolor[rgb]{0.9,0.9,0.9}{
            \makecell[{{p{\linewidth}}}]{
                \texttt{\tiny{[GM$|$GM]}}
                \texttt{continue} \\
            }
        }
    }
    & & \\ \\

    \theutterance \stepcounter{utterance}  
    & & & \multicolumn{2}{p{0.3\linewidth}}{
        \cellcolor[rgb]{0.9,0.9,0.9}{
            \makecell[{{p{\linewidth}}}]{
                \texttt{\tiny{[GM$|$GM]}}
                \texttt{success} \\
            }
        }
    }
    & & \\ \\

\end{supertabular}
}

\end{document}
